% ============================================================
% MOBILE AND WEB DESIGN
% Responsible member: Nguyen Dao Quoc Khanh
% ============================================================

\section{Mobile and Web Design}

This section presents the proposed user interfaces for the Photographer Mobile
Application, Film Lab Web Portal, Administration Web Portal, Marketplace, and
Digital Film Archive. Flutter is proposed for the mobile application, while
ReactJS or Next.js is proposed for the web portals.

\subsection{UI/UX Guidelines}

The interfaces follow these common guidelines:

\begin{itemize}
    \item The mobile application uses simple bottom navigation for Home,
    Orders, Archive, Marketplace, and Profile.
    \item The Film Lab and Administration portals use a consistent sidebar,
    page header, searchable tables, filters, and pagination.
    \item Each screen has one clear primary action, such as booking a service,
    updating an order status, or uploading scan files.
    \item Order states use both color and text labels. Important examples are
    Pending, Processing, Scanning, Ready for Delivery, and Completed.
    \item Forms show visible labels and short validation messages without
    removing valid information entered by the user.
    \item Loading, empty, success, confirmation, and error states are provided
    for important actions.
    \item Destructive actions and actions that affect another user require a
    confirmation step.
    \item Mobile and web layouts remain readable on their supported screen
    sizes and provide keyboard focus for web controls.
\end{itemize}

\subsection{Photographer Mobile Application}

The Photographer Application is developed with Flutter for Android and iOS.
It allows photographers to find a Film Lab, create processing orders, track
their film, receive scan files, use the Marketplace, and manage a personal
Digital Film Archive.

\begin{table}[htbp]
\centering
\small
\renewcommand{\arraystretch}{1.1}
\caption{Main Photographer Mobile Application Screens}
\label{tab:photographer-mobile-screens}
\begin{tabular}{|p{3.4cm}|p{9.2cm}|}
\hline
\textbf{Screen Group} & \textbf{Main Functions} \\
\hline
Authentication & Register, log in, reset password, and verify an account. \\
\hline
Home & Show current orders, nearby Film Labs, recommendations, and
notifications. \\
\hline
Film Lab Discovery & Search and filter by location, film format, service,
price, rating, and processing time. \\
\hline
Booking & Select a service, enter film information, choose pickup or drop-off,
review the cost, and confirm the order. \\
\hline
Orders & Display order history and the processing status timeline. \\
\hline
Scan Delivery & Preview thumbnails, download scan files, and save photographs
to the archive. \\
\hline
Marketplace and Community & View listings and community content, contact
sellers, and report inappropriate content. \\
\hline
AI Assistant & Ask questions about film stocks, exposure, camera settings,
Film Labs, and processing services. \\
\hline
Profile & Update personal details, saved addresses, notifications, and account
settings. \\
\hline
\end{tabular}
\end{table}

The main booking flow is:

\begin{center}
\texttt{Login $\rightarrow$ Search Film Lab $\rightarrow$ Compare Services
$\rightarrow$ Enter Film Details $\rightarrow$ Choose Pickup
$\rightarrow$ Confirm Order $\rightarrow$ Track Status}
\end{center}

Before confirmation, the application displays the service price, pickup fee,
estimated completion time, and total amount. If location permission is not
available, the photographer can manually enter a district or city.

\subsection{Film Lab Web Portal}

The Film Lab Web Portal gives owners and employees a larger workspace for
handling orders and scan files. A Lab owner can manage the Film Lab profile,
services, prices, employees, and reports. A Lab employee can process assigned
orders and upload scan files.

The main pages are Dashboard, Orders, Order Details, Processing Board,
Schedule, Scan Delivery, Services and Pricing, Customers, Reports, and Lab
Settings. The processing workflow is:

\begin{center}
\texttt{New Order $\rightarrow$ Accepted $\rightarrow$ Film Received
$\rightarrow$ Developing $\rightarrow$ Scanning $\rightarrow$
Quality Checking $\rightarrow$ Scans Published $\rightarrow$ Completed}
\end{center}

Every status update records the employee, update time, and an optional note.
Invalid status jumps are disabled. During scan delivery, the portal validates
file type and size, displays upload progress, supports retry, and requires a
final confirmation before publishing files to the photographer.

\subsection{Administration Web Portal}

The Administration Portal supports system administrators, moderators, finance
staff, and support staff. Role-based permissions limit each user to the data
and actions required for their work.

The portal provides the following main functions:

\begin{itemize}
    \item manage users, roles, account status, and Film Lab applications;
    \item manage service categories, platform fees, and public settings;
    \item monitor payments, failed transactions, refunds, and activity;
    \item review community content, marketplace listings, complaints, and
    disputes;
    \item view reports and search the administrative audit log.
\end{itemize}

For Film Lab approval, an administrator checks the submitted business and
service information, requests missing information when necessary, and records
an approval or rejection reason. Important access, moderation, payment, and
configuration actions are written to the audit log.

\subsection{Marketplace Interface}

The Marketplace supports cameras, lenses, film rolls, accessories, and other
film photography equipment. Users can search by keyword and filter by category,
brand, condition, location, price, and listing date.

The main screens are Marketplace Home, Search Results, Listing Details, Create
Listing, Seller Profile, Messages, Saved Items, My Listings, and Transactions.
A seller creates a listing using the following flow:

\begin{center}
\texttt{Select Category $\rightarrow$ Upload Photos $\rightarrow$
Enter Item Details $\rightarrow$ Set Price and Delivery
$\rightarrow$ Preview $\rightarrow$ Submit}
\end{center}

The interface clearly marks Draft, Active, Reserved, Sold, Hidden, Rejected,
and Expired listings. Users can report suspicious listings, messages, or
accounts, and the report is sent to the Administration Portal.

\subsection{Digital Film Archive Interface}

The Digital Film Archive stores scan files published by Film Labs and helps the
photographer organize them by order, album, film stock, shooting date, tags,
and favorite status.

The archive includes a library, order scan view, photo viewer, album manager,
search and filters, metadata editor, downloads, and storage and privacy
settings. Film Lab metadata, such as order code, scanner, resolution, and
delivery date, is separated from personal metadata entered by the photographer.
Changing an album, tag, or note does not overwrite the original delivered file.

New scan deliveries are private by default. Sharing requires a separate user
action. Large downloads show progress and can be retried if the network is
interrupted.

\subsection{Frontend Component Architecture}

The Flutter application and ReactJS or Next.js portals use the same separation
of responsibilities.

\begin{table}[htbp]
\centering
\small
\renewcommand{\arraystretch}{1.1}
\caption{Frontend Architecture Layers}
\label{tab:frontend-architecture-layers}
\begin{tabular}{|p{3.2cm}|p{9.4cm}|}
\hline
\textbf{Layer} & \textbf{Responsibility} \\
\hline
Presentation & Screens, pages, layouts, forms, tables, dialogs, and reusable
interface components. \\
\hline
State and Feature & Authentication, order, marketplace, archive, profile, and
notification state. \\
\hline
Domain & Shared models and application rules for users, Film Labs, services,
orders, scans, and listings. \\
\hline
API and Service & REST API communication, file upload and download, and
real-time updates. \\
\hline
Shared & Theme, routing, form validation, error mapping, localization, and
common utilities. \\
\hline
\end{tabular}
\end{table}

A user action starts in the presentation layer. The feature state validates the
input and calls the relevant service. The service communicates with the backend
or file storage and returns a shared model. Finally, the state is updated and
the interface shows the result. Screens do not store secrets or call backend
endpoints directly.

\subsection{Main Cross-Platform UI Flow}

The main interaction between the mobile and web interfaces is summarized as
follows:

\begin{enumerate}
    \item The photographer searches for a Film Lab and creates an order in the
    mobile application.
    \item Film Lab staff review the order and update its processing stages in
    the web portal.
    \item Status notifications are displayed in the Photographer Application.
    \item The Film Lab uploads and publishes scan files after quality checking.
    \item The photographer previews, downloads, and organizes the scans in the
    Digital Film Archive.
    \item Reports or disputes are transferred to the Administration Portal for
    review.
\end{enumerate}

\subsection{Frontend Sequence Summaries}

Four frontend sequence diagrams are planned for implementation and
presentation:

\begin{enumerate}
    \item \textbf{Create service order:} the Photographer App validates the
    booking, sends it to the Platform API, and displays the returned order code.
    \item \textbf{Update processing status:} the Film Lab Portal sends the next
    status to the API, which records the history and notifies the photographer.
    \item \textbf{Publish scans:} the Film Lab Portal uploads files to storage,
    saves metadata through the API, and makes the files available in the mobile
    application.
    \item \textbf{Moderate a listing:} the Administration Portal retrieves the
    report and evidence, records the decision, updates the Marketplace, and
    writes an audit event.
\end{enumerate}

\subsection{Initial Wireframes and Presentation Plan}

The first wireframes should show Photographer Home, Film Lab Search, Film Lab
Details, Booking, Order Tracking, Scan Delivery, Film Lab Dashboard, Processing
Board, Scan Upload, Admin Moderation, Marketplace Search, Listing Details,
Archive Library, and Photo Viewer.

During the presentation, the UI flow should be introduced first, followed by
one mobile example, one web portal example, the component architecture, and the
scan delivery sequence. The recommended speaking time for this section is four
to five minutes.
